%% LyX 2.3.6.1 created this file.  For more info, see http://www.lyx.org/.
%% Do not edit unless you really know what you are doing.
\documentclass{article}
\usepackage[utf8]{inputenc}
\usepackage{amsmath}
\usepackage{amssymb}

\makeatletter

%%%%%%%%%%%%%%%%%%%%%%%%%%%%%% LyX specific LaTeX commands.
\DeclareTextSymbolDefault{\textquotedbl}{T1}

%%%%%%%%%%%%%%%%%%%%%%%%%%%%%% User specified LaTeX commands.
\usepackage{amsfonts}
\usepackage{algorithm}
\usepackage{algpseudocode}

\makeatother

\begin{document}

\section{Introdução e Motivação}

A motivação central deste trabalho reside no fato de que o preço não
é uma variável física independente, mas o resultado da interação coletiva
de agentes (traders, algoritmos, investidores institucionais). 

Existe toda uma fenomenologia baseada em padrões:
\begin{itemize}
\item Traders técnicos são treinados para identificar \textquotedbl padrões\textquotedbl{}
(ombro-cabeça-ombro, suportes, resistências, gaps).
\item Modelos quantitativos utilizados por investidores instucionais calculam
\textquotedbl fatores\textquotedbl{} como volatilidade ou momentum
para estruturar carteiras de investimento.
\item Algoritmos como HFTs (High Frequency Trading) reagem a pequenas variações
de preço de forma matemática, buscando explorar ineficiências de curtíssimo
prazo. 
\item Volatilidade: a volatilidade dos ativos financeiros é inconstante
ao longo do tempo. Em períodos de crise ou instabilidade econômica,
essa volatilidade cresce, ao mesmo tempo em que os preços das ações
caem. A volatilidade também forma clusters, isto é, períodos de volatilidade
alta ou baixa tendem a se agrupar, indicando uma persistência temporal.
\item Mudanças de regime: é comum classificar os mercados como de alta ou
de baixa (Bull/Bear Market), mercado lateralizado, etc, indicando
movimentos dominantes.
\end{itemize}
Essas observações empíricas mostram que séries de preços financeiros
exibem simultaneamente aleatoriedade e padrões recorrentes. A presença
de estruturas estatísticas persistentes sugere que os preços não evoluem
puramente de forma randômica, mas sim como combinações dinâmicas de
padrões elementares.

Este trabalho propõe que tais séries são projeções observáveis de
uma dinâmica estocástica ocorrendo em um espaço de padrões estruturais,
aqui chamados de \emph{primitivas}. Essas \emph{primitivas} são autovetores
de um operador integral $K$\cite{Baaquie2004}, que possui uma dinâmica
própria a qual impacta diretamente na dinâmica dos preços observados.
As \emph{primitivas} emergem como modos estruturais desse operador,
e podem ser estimadas empiricamente via PCA (Principal Component Analysis).
O conjuto de \emph{primitivas} define um \emph{regime} do mercado,
e este \emph{regime} é definido pela geometria do operador $K$.

\section{Formulação básica}

Considere uma janela temporal finita

\[
t\in[0,T],
\]
onde $t$ é o índice interno da janela. Definimos os retornos $P(t)$
a partir do preço $S(t)$ do ativo como

\[
P(t)=\ln\left(\frac{S(t)}{S(t-\Delta t)}\right).
\]
A covariância temporal dentro da janela é dada por

\[
K^{(r)}(t,t')=\langle P(t),P(t')\rangle^{(r)},
\]
onde a média é feita sobre uma coleção de janelas onde o regime $r$
é o mesmo. Esse é um kernel simétrico e definido positivo em $L^{2}([0,T])$. 

O PCA resolve a equação integral

\[
\int_{0}^{T}K^{(r)}(t,t')\phi_{k}^{(r)}(t')dt'=\lambda_{k}^{(r)}\phi_{k}^{(r)}(t),
\]
onde as funções $\phi_{k}^{(r)}(t)$ formam uma base ortonormal em
$L^{2}([0,T])$

\begin{equation}
\int_{0}^{T}\phi_{k}^{(r)}(t)\phi_{j}^{(r)}(t)dt=\delta_{kj}.\label{eq:ortonormalidade}
\end{equation}
Dentro da janela, os coeficientes não dependem do tempo,de forma que
podemos escrever $P(t)$ como
\[
P(t)=\sum_{k=1}^{m}a_{k}^{(r)}\phi_{k}^{(r)}(t)+\eta^{(r)}(t),
\]
onde

\[
a_{k}^{(r)}=\int_{0}^{T}P(t)\phi_{k}^{(r)}(t),dt
\]
são constantes para aquela janela. O resíduo é

\[
\eta^{(r)}(t)=\sum_{k=m+1}^{\infty}a_{k}^{(r)}\phi_{k}^{(r)}(t).
\]

Ao passar para a próxima janela, temos duas possibilidades. Se o regime
permanece o mesmo, o kernel é o mesmo, e terá os mesmos autovalores
$\lambda_{k}^{(r)}$ e autovetores $\phi_{k}^{(r)}(t)$, mas o novo
histórico $P(t)$ gera novos coeficientes $a_{k}^{(r)}\longrightarrow\tilde{a}_{k}^{(r)}$.
Por outro lado, se há uma mudança de regime, $r\to r'$, o kernel
muda $K^{(r)}(t,t')\longrightarrow K^{(r')}(t,t')$, e mudam as autofunções
$\phi_{k}^{(r)}(t)$, os autovalores $\lambda_{k}^{(r)}$ e o resíduo
$\eta^{(r)}(t)$.

Dentro de um regime $r$, o modelo não tem nenhuma dependência temporal
adicional. O único objeto que realmente carrega a informação do regime
é o kernel $K^{(r)}(t,t')$, logo a mudança de regime é a mudança
deste kernel. Enquanto modelos econométricos tradicionais tratam mudanças
de regime via cadeias de Markov discretas\cite{Hamilton1989}, o presente
modelo propõe uma transição via dinâmica espectral do operador de
covariância.

\section{Dinâmica do kernel $K$}

As autofunções satisfazem, para cada $k$,

\begin{equation}
\int_{0}^{T}K(t,t';\tau)\phi_{k}(t';\tau)dt'=\lambda_{k}(\tau)\phi_{k}(t;\tau),\label{eq:autovalor_K}
\end{equation}
onde introduzimos explicitamente $\tau$, o índice da janela, que
é o tempo lento, não o tempo interno $t\in[0,T]$. É em $\tau$ que
o regime pode mudar. Suponha agora que o regime varia lentamente,
ou seja, $K(t,t';\tau)$ é diferenciável em $\tau$. Derivamos (\ref{eq:autovalor_K})
em relação a $\tau$

\[
\int_{0}^{T}\frac{\partial K}{\partial\tau}(t,t';\tau)\phi_{k}(t';\tau)dt'+\int_{0}^{T}K(t,t';\tau)\frac{\partial\phi_{k}}{\partial\tau}(t';\tau)dt'=\frac{d\lambda_{k}}{d\tau}\phi_{k}(t;\tau)+\lambda_{k}(\tau)\frac{\partial\phi_{k}}{\partial\tau}(t;\tau),
\]
como $\{\phi_{j}\}$ forma uma base ortonormal do espaço $L^{2}$,
multiplicamos a equação acima por $\phi_{j}(t)$ e integramos em relação
a $t$ no intervalo $[0,T]$, obtendo

\[
\begin{aligned}\int_{0}^{T}\phi_{j}(t)\left[\int_{0}^{T}\frac{\partial K}{\partial\tau}(t,t')\phi_{k}(t')dt'\right]dt & +\int_{0}^{T}\phi_{j}(t)\left[\int_{0}^{T}K(t,t')\frac{\partial\phi_{k}}{\partial\tau}(t')dt'\right]dt\\
 & =\int_{0}^{T}\phi_{j}(t)\left[\frac{\partial\lambda_{k}}{\partial\tau}\phi_{k}(t)+\lambda_{k}\frac{\partial\phi_{k}}{\partial\tau}(t)\right]dt.
\end{aligned}
\]
Como $K$ é autoadjunto, podemos inverter a ordem de integração no
segundo termo do lado esquerdo desta equação

\[
\int_{0}^{T}\phi_{j}(t)\int_{0}^{T}K(t,t')\frac{\partial\phi_{k}}{\partial\tau}(t')dt'dt=\int_{0}^{T}\frac{\partial\phi_{k}}{\partial\tau}(t')\left[\int_{0}^{T}K(t',t)\phi_{j}(t)dt\right]dt'
\]
Usando a equação espectral (\ref{eq:autovalor_K}), o termo torna-se

\[
\lambda_{j}\int_{0}^{T}\phi_{j}(t)\frac{\partial\phi_{k}}{\partial\tau}(t)dt.
\]
No lado direito, usamos a ortonormalidade na equação (\ref{eq:ortonormalidade}),
e obtemos para $j\neq k$

\[
\begin{aligned}\int_{0}^{T}\int_{0}^{T}\phi_{j}(t)\frac{\partial K}{\partial\tau}(t,t')\phi_{k}(t')dt'dt+\lambda_{j}\int_{0}^{T}\phi_{j}(t)\frac{\partial\phi_{k}}{\partial\tau}(t)dt & =\lambda_{k}\int_{0}^{T}\phi_{j}(t)\frac{\partial\phi_{k}}{\partial\tau}(t)dt\\
\int_{0}^{T}\int_{0}^{T}\phi_{j}(t)\frac{\partial K}{\partial\tau}(t,t')\phi_{k}(t')dt'dt & =(\lambda_{k}-\lambda_{j})\int_{0}^{T}\phi_{j}(t)\frac{\partial\phi_{k}}{\partial\tau}(t)dt
\end{aligned}
,
\]
que pode ser reescrita como
\begin{equation}
\int_{0}^{T}\phi_{j}(t)\frac{\partial\phi_{k}}{\partial\tau}(t)dt=\frac{{\displaystyle \int_{0}^{T}\int_{0}^{T}\phi_{j}(t)\frac{\partial K}{\partial\tau}(t,t')\phi_{k}(t')dtdt'}}{\lambda_{k}-\lambda_{j}}.\label{eq:adiabatica0}
\end{equation}
Como $\{\phi_{j}\}$ é uma base completa, qualquer função no espaço,
incluindo $\frac{\partial\phi_{k}}{\partial\tau}$, pode ser escrita
como

\[
\frac{\partial\phi_{k}}{\partial\tau}(t)=\sum_{j}c_{jk}\phi_{j}(t),
\]
onde os coeficientes $c_{jk}$ são dados pela projeção

\[
c_{jk}=\int_{0}^{T}\phi_{j}(t)\frac{\partial\phi_{k}}{\partial\tau}(t)dt.
\]
Para $j=k$, a diferenciação da condição de normalização $\int\phi_{k}^{2}dt=1$
resulta em

\[
\frac{\partial}{\partial\tau}\int_{0}^{T}\phi_{k}^{2}(t)dt=2\int_{0}^{T}\phi_{k}(t)\frac{\partial\phi_{k}}{\partial\tau}(t)dt=0\implies c_{kk}=0.
\]
Multiplicando a equação (\ref{eq:adiabatica0}) por $\phi_{j}$e somando
para $j\neq k$ na expansão em série, obtemos finalmente

\begin{equation}
\frac{\partial\phi_{k}}{\partial\tau}(t)=\sum_{j\neq k}\frac{{\displaystyle \int\int\phi_{j}(t_{1})\frac{\partial K}{\partial\tau}(t_{1},t_{2})\phi_{k}(t_{2})dt_{1}dt_{2}}}{\lambda_{k}-\lambda_{j}}\phi_{j}(t),\label{eq:adiabatica}
\end{equation}

que é a \emph{dinâmica adiabática das bases do PCA} neste contexto.
Esta equação é análoga à teoria de perturbação adiabática de operadores
auto-adjuntos \cite{Messiah1961,Kato1995}. Se $K$ muda lentamente,
as $\phi_{k}$ giram lentamente no espaço $L^{2}$. Se dois autovalores
se aproximam $\lambda_{k}\approx\lambda_{j}$, a base gira violentamente,
estes autovetores tornam-se linearmente dependentes, e há uma mudança
em todo o espectro de autovalores e autovetores, o que representa
uma mudança de mudança de regime. Observe que esta dinâmica adiabática
de bases, Eq.(\ref{eq:adiabatica}), só se aplica aos modos fora do
bulk de MP, pois as bases dentro do bulk são essencialmente ruído.

Suponha agora que $K$ também se movimenta por saltos. Assim em $\tau=\tau_{0}$

\begin{equation}
K(\tau_{0}^{+})=K(\tau_{0}^{-})+\Delta K.\label{eq:saltos}
\end{equation}
Aqui não existe derivada, e a equação (\ref{eq:adiabatica}) não vale.
O que ocorre é uma rotação instantânea da base. As novas autofunções
são combinações das antigas

\[
\phi_{k}^{+}(t)=\sum_{j}M_{kj}\phi_{j}^{-}(t),
\]
onde

\[
M_{kj}=\int_{0}^{T}\phi_{k}^{+}(t)\phi_{j}^{-}(t)dt
\]
é uma matriz ortogonal. Isso é uma transição não-adiabática. 

Propomos então um modelo com difusão estocástica lenta

\[
dK=A[K]d\tau+B[K]dW
\]
para modelar a parte adiabática de K. Neste caso temos rotação irregular
dos autovetores, e variações realistas dentro do regime, com a possibilidade
de mudanças estruturais através do cruzamento de autovalores. De forma
geral, a dinâmica de $K$ pode ser modelada como a superposição destas
contribuições, uma parte lenta com difusão estocástica $\frac{\partial K}{\partial\tau}=A[K]+B[K]dW$
que gera a equação diferencial (\ref{eq:adiabatica}), e uma parte
de saltos dada pela equação (\ref{eq:saltos}). A forma mais geral
é então um processo de Itô com saltos no espaço de kernels
\begin{equation}
dK=A[K]d\tau+B[K]dW_{\tau}+\Delta KdN_{\tau},\label{eq:ito}
\end{equation}
onde $A[K]$ é um termo de drift determinístico, $B[K]$ é a intensidade
do ruído, $dW$ é um processo de Wiener, e $dN$ é um processo de
Poisson, relacionado a mudanças violentas de regime em grandes crises,
como a do vírus COVID.

\section{Dinâmica dos coeficientes $a_{k}$}

Dentro da janela

\[
P(t)=\sum_{k}a_{k}(\tau)\phi_{k}(t;\tau),
\]
com

\begin{equation}
a_{k}(\tau)=\int_{0}^{T}P(t)\phi_{k}(t;\tau)dt.\label{eq:aktempo}
\end{equation}
Como $a_{k}$ muda quando a base gira , derivamos (\ref{eq:aktempo})
em relação a $\tau$

\[
\frac{da_{k}}{d\tau}=\int_{0}^{T}P(t)\frac{\partial\phi_{k}}{\partial\tau}(t)dt,
\]
substituindo (\ref{eq:adiabatica}) obtemos

\[
\frac{da_{k}}{d\tau}=\sum_{j\neq k}\frac{{\displaystyle \int\int\phi_{j}(t_{1})\frac{\partial K}{\partial\tau}(t_{1},t_{2})\phi_{k}(t_{2})dt_{1}dt_{2}}}{\lambda_{k}-\lambda_{j}}\underbrace{\int_{0}^{T}P(t)\phi_{j}(t)dt}_{a_{j}}.
\]
Logo,
\[
\frac{da_{k}}{d\tau}=\sum_{j\neq k}\Omega_{kj}(\tau)a_{j}(\tau),
\]
onde
\[
\Omega_{kj}(\tau)=\frac{{\displaystyle \int_{0}^{T}\int_{0}^{T}\phi_{j}(t_{1})\frac{\partial K}{\partial\tau}(t_{1},t_{2})\phi_{k}(t_{2})dt_{1}dt_{2}}}{\lambda_{k}-\lambda_{j}},
\]
ou seja, os coeficientes não são independentes, eles obedecem um sistema
acoplado. Quando o kernel gira, a energia de um modo passa para outros,
o que se reflete em um aumento do resíduo e a perda de estabilidade
do PCA antes da mudança de regime. Por outro lado, quando ocorre um
salto $dN_{\tau}=1$,

\[
K^{+}=K^{-}+\Delta K,
\]
as autofunções mudam

\[
\phi_{k}^{+}=\sum_{j}M_{kj}\phi_{j}^{-},
\]
e os coeficientes sofrem rotação instantânea
\[
a_{k}^{+}=\sum_{j}M_{kj}a_{j}^{-}.
\]


\section{Ruído dos resíduos}

O ruído $\eta$ não é arbitrário, ele é o ruído previsto pela teoria
de matrizes aleatórias (RMT) para o kernel temporal. Da decomposição
espectral completa do kernel em um regime $r$,

\[
K(t,t')=\sum_{k=1}^{\infty}\lambda_{k}\phi_{k}(t)\phi_{k}(t')
\]

separamos os modos estruturais $\lambda_{k}>\lambda^{+}$e os modos
de ruído (bulk) $\lambda_{k}\le\lambda^{+}$, onde 
\[
\lambda^{+}=\sigma^{2}(1+\sqrt{q})^{2}
\]
é o limite superior de Marchenko--Pastur (MP) \cite{Laloux1999,Plerou1999},
com $q=L/N$ (tamanho da janela / número de amostras). A separação
entre componentes informativos e o ruído universal de matrizes aleatórias
é validada empiricamente para grandes conjuntos de ativos\cite{Plerou2002}.
A reconstrução truncada usa $m$ modos informativos

\[
P(t)=\sum_{k=1}^{m}a_{k}\phi_{k}(t)+\eta(t),
\]
com
\[
\eta(t)=\sum_{k>m}a_{k}\phi_{k}(t).
\]
Como a base diagonaliza a covariância,

\[
\langle a_{k},a_{j}\rangle=\lambda_{k}\delta_{kj}
\]
Para $k>m$, temos

\[
\lambda_{k}\in\text{espectro MP}.
\]
De acordo com MP, esses autovalores são aqueles de uma matriz de covariância
de ruído branco finito-amostral, ou seja,

\[
a_{k}\sim\mathcal{N}(0,\lambda_{k})\quad\text{independentes}.
\]

Agora calculamos a correlação do bulk

\begin{equation}
\langle\eta(t),\eta(t')\rangle=\sum_{k>m}\lambda_{k}\phi_{k}(t)\phi_{k}(t'),\label{eq:correruido0}
\end{equation}
que é a parte do kernel associada ao ruído
\[
K_{\text{ruído}}(t,t')=\sum_{\lambda_{k}\le\lambda_{+}}\lambda_{k}\phi_{k}(t)\phi_{k}(t').
\]
Segundo MP, para ruído branco de variância $\sigma^{2}$, o espectro
da covariância empírica ocupa o intervalo

\[
\lambda\in\left[\sigma^{2}(1-\sqrt{q})^{2},\sigma^{2}(1+\sqrt{q})^{2}\right],
\]
logo, os modos da cauda são modos que não carregam estrutura temporal
real, são artefatos estatísticos. Da eq. (\ref{eq:correruido0}) temos
a variância pontual do bulk

\[
\langle\eta(t)^{2}\rangle=\sum_{k>m}\lambda_{k}\phi_{k}(t)^{2}
\]
e a sua correlação temporal

\[
\langle\eta(t),\eta(t')\rangle=K_{\text{ruído}}(t,t'),
\]
mas como os autovetores$\phi_{k}$ da cauda são essencialmente aleatórios
\cite{CouilletDebbah2011,Paul2007} e não correlacionados, essa soma
converge para
\[
\langle\eta(t),\eta(t')\rangle\approx\sigma_{\eta}^{2}\delta(t-t'),
\]
ou seja, o resíduo é aproximadamente ruído branco temporal dentro
da janela. Note que apesar de termos um ruído residual branco, isso
não significa que os retornos sejam não correlacionados, uma parte
da correlação é devida aos modos fora do bulk, o que pode resultar
em correlações de longo alcance (cauda longa) \cite{Mantegna1999}.
Para obtermos a variância total do resíduo $\sigma_{\eta}^{2}$, integramos
$\langle\eta(t)^{2}\rangle$sobre o domínio temporal $T$ e usamos
a normalização das autofunções 
\[
\begin{aligned}\sigma_{\eta}^{2} & =\int_{0}^{T}\langle\eta(t)^{2}\rangle dt\\
 & =\int_{0}^{T}\sum_{k>m}\lambda_{k}\phi_{k}(t)^{2}dt\\
 & =\sum_{k>m}\lambda_{k}\int_{0}^{T}\phi_{k}(t)^{2}dt
\end{aligned}
,
\]
e como $\int\phi_{k}(t)^{2}dt=1$, obtemos o resultado final para
a variância acumulada do ruído

\[
\sigma_{\eta}^{2}=\sum_{k>m}\lambda_{k}.
\]
Esta variância aumenta quando os autovalores informativos cruzam o
limiar de MP dado por $\lambda^{+}$, o que sinaliza uma mudança de
regime.

\section{Mudanças de regime e transições de fase}

Enquanto os autovalores são bem separados,$\lambda_{k}\gg\lambda_{j}$,
e a equação que obtivemos

\[
\frac{\partial\phi_{k}}{\partial\tau}\propto\frac{1}{\lambda_{k}-\lambda_{j}}
\]
implica que a base é rígida e cada $\phi_{k}$ tem identidade própria,
ou seja, o PCA é estável. Isso é o que se observa dentro de um regime.
Por outro lado quando $\lambda_{k}\to\lambda_{j}$, o denominador
vai a zero, ou seja, a derivada de $\phi_{k}$ tende ao infinito$\frac{\partial\phi_{k}}{\partial\tau}\to\infty$.
Com isto as autofunções deixam de ter identidade individual e passam
a poder se misturar arbitrariamente dentro do subespaço gerado por
elas. Matematicamente, quando há degenerescência espectral, os autovetores
não são mais únicos, só o subespaço é bem definido. Antes de uma mudança
de regime os primeiros modos começam a “trocar de forma”, ficam instáveis,
a reconstrução piora e o resíduo cresce. Não é somennte uma mudança
em um autovetor. A estrutura espectral inteira muda, porque quando
dois modos se tornam quase degenerados, o subespaço relevante muda
de orientação, os coeficientes $a_{k}$ passam a trocar energia fortemente
e a separação entre modo 1, modo 2, etc, deixa de fazer sentido, o
que força uma reorganização global da base.

Em teoria espectral e física estatística, uma transição de fase ocorre
quando uma pequena variação de um parâmetro produz mudança não suave
na estrutura espectral do operador. Aqui o parâmetro é $\tau$ (tempo
das janelas), e a não-suavidade ocorre quando$\lambda_{k}-\lambda_{j}\to0,$ou
seja, o mercado passa por transições de fase espectrais do operador
de covariância temporal, onde os modos deixam de ser bem definidos,
o ruído aumenta, a base inteira se reorganiza e surge um novo conjunto
estável depois. O parâmetro de ordem desta transição é 
\[
\Delta(\tau)=min_{k\neq j}|\lambda_{k}-\lambda_{j}|,
\]
ou seja, quando $\Delta\rightarrow0$ temos a transição. Da mesma
forma que antes, observe que essa transição de fase só faz sentido
para os modos fora do bulk de MP. Esse processo de degenerescência
também pode fazer com que um dos autovalores estruturais cruze o limite
de MP, resultado de uma perda de estrutura real do mercado, ou seja,
a dinâmica se aproxima da dinâmica desestruturada do bulk, que é um
ruído branco. Por outro lado, pode haver o movimento inverso, quando
há estruturação do mercado, um dos autovalores que estavam no bulk
pode cruzar o limiar de MP. Assim, temos dois mecanismos principais
de mudança de regime, um onde a base gira violentamente e se reorganiza,
mantendo o número de autovalores estruturais, e outra onde há mudança
na quantidade de autovalores fora do bulk, que representa mudanças
na quantidade de informação disponível. 

\section{Teorma da Persistência Estrutural}

Nesta seção estabelecemos formalmente que os autovetores do operador
de covariância empírico associados a autovalores fora do bulk de Marchenko--Pastur
são estatisticamente consistentes e convergem para os modos próprios
do processo gerador subjacente. Esta propriedade justifica rigorosamente
a interpretação de que as primitivas temporais observadas via PCA
móvel correspondem a estruturas dinâmicas reais do mercado, e não
a artefatos de amostragem finita.

Considere que, dentro de um regime, os retornos normalizados possam
ser representados por

\begin{equation}
P(t)=\sum_{k=1}^{m}a_{k}\,\psi_{k}(t)+\epsilon(t),
\end{equation}
onde $\psi_{k}(t)\in L^{2}([0,T])$ são funções determinísticas representando
componentes estruturais de baixa dimensão, $a_{k}$ são variáveis
aleatórias com variância finita e $\epsilon(t)$ é ruído branco temporal
com variância $\sigma^{2}$. O operador de covariância real do processo
é então

\begin{equation}
K_{\text{real}}(t,t')=\sum_{k=1}^{m}\mathbb{E}[a_{k}^{2}]\,\psi_{k}(t)\psi_{k}(t')+\sigma^{2}\delta(t-t').
\end{equation}
Esse operador é a soma de um operador de posto finito com um operador
identidade escalado pelo ruído branco, possuindo portanto autovalores
isolados associados às funções $\psi_{k}$. Na prática, estimamos
a covariância a partir de $N$ janelas pertencentes ao mesmo regime

\begin{equation}
K_{\text{emp}}(t,t')=\frac{1}{N}\sum_{i=1}^{N}P_{i}(t)P_{i}(t').
\end{equation}

Resultados clássicos de estatística funcional garantem que, quando
$N\to\infty$, $\|K_{\text{emp}}-K_{\text{real}}\|_{\mathrm{HS}}\to0,$onde
$\|\cdot\|_{\mathrm{HS}}$ é a norma de Hilbert--Schmidt \cite{Bosq2000}.
O operador $K_{\text{real}}$ possui $m$ autovalores destacados acima
do contínuo associado ao ruído branco. A teoria de matrizes aleatórias
para modelos spiked mostra que, quando existe ruído de alta dimensão
com razão dimensional fixa $q=T/N$, os autovalores empíricos associados
às componentes estruturais se separam do bulk de Marchenko--Pastur
sempre que os autovalores reais excedem um limiar crítico \cite{Johnstone2001,Baik2005}.
Além disso, os autovetores empíricos correspondentes apresentam sobreposição
finita com os modos reais, isto é, $|\langle\phi_{k}^{(\text{emp})},\psi_{k}\rangle|^{2}\to c_{k}>0.$
Como $K_{\text{emp}}\to K_{\text{real}}$ e os autovalores estruturais
são isolados, podemos aplicar a teoria de perturbação espectral para
operadores auto-adjuntos compactos \cite{Kato1995,DavisKahan1970}.
Isso implica que os autovetores empíricos convergem em norma para
os autovetores reais, ouu seja, $\|\phi_{k}^{(\text{emp})}-\psi_{k}\|\to0.$
Assim, podemos enunciar o \textbf{teorema da persistência estrutural:} 

\emph{Sob o modelo acima, o operador de covariância empírico converge
para o operador real em norma de Hilbert--Schmidt. Os autovalores
isolados de $K_{\text{real}}$ geram autovalores empíricos que se
separam do bulk de Marchenko--Pastur, e os autovetores correspondentes
convergem para os modos estruturais $\psi_{k}$ quando $N,T\to\infty$
com razão dimensional fixa.}

Esse resultado garante que as primitivas temporais identificadas via
PCA não são artefatos estatísticos, mas correspondem a graus de liberdade
dinâmicos reais do processo gerador do mercado. 

\section{Dinâmica Estocástica do Operador $K$ e Emergência do Bulk Espectral}

Assumimos que o operador estrutural $K(\tau)$ evolui segundo uma
dinâmica de Itô no tempo lento $\tau$ \cite{ito1944,dynkin1965}

\begin{equation}
dK=A(K,\tau)\,d\tau+d\eta,
\end{equation}
onde $A$ é um operador determinístico (não necessariamente linear)
e $d\eta$ é um incremento estocástico com

\begin{equation}
\mathbb{E}[d\eta]=0,\qquad\mathbb{E}[d\eta^{ab}d\eta^{cd}]=B^{ab,cd}\,d\tau.
\end{equation}
Note que não assumimos nenhuma forma específica para $A$ ou para
$B$, apenas que os momentos de segunda ordem existam. A solução integral
da equação acima pode ser escrita formalmente como

\begin{equation}
K(\tau)=\Phi(\tau,0)K(0)+\int_{0}^{\tau}\Phi(\tau,s)\,d\eta(s),
\end{equation}
onde $\Phi(\tau,s)$ é o propagador determinístico associado à parte
$A$. Logo, toda a aleatoriedade de $K$ provém do termo integral
estocástico. Sejam $(\lambda_{k}(\tau),\phi_{k}(\tau))$ os autovalores
e autovetores instantâneos de $K$

\begin{equation}
K(\tau)\phi_{k}(\tau)=\lambda_{k}(\tau)\phi_{k}(\tau),
\end{equation}
Aplicando cálculo de Itô matricial e teoria de perturbação estocástica
de autovalores \cite{dyson1962,anderson2010}, obtemos

\begin{align}
d\lambda_{k} & =\langle\phi_{k},dK\,\phi_{k}\rangle+\sum_{j\neq k}\frac{|\langle\phi_{j},dK\,\phi_{k}\rangle|^{2}}{\lambda_{k}-\lambda_{j}},\\
d\phi_{k} & =\sum_{j\neq k}\frac{\langle\phi_{j},dK\,\phi_{k}\rangle}{\lambda_{k}-\lambda_{j}}\phi_{j}.
\end{align}
Assim, o termo determinístico afeta suavemente o espectro e o termo
estocástico introduz mistura espectral. Projetando o ruído nos autovetores
de $K$

\begin{equation}
\Xi_{jk}=\langle\phi_{j},d\eta\phi_{k}\rangle,
\end{equation}
a variância do incremento dos autovetores é

\begin{equation}
\mathbb{E}\big[\|d\phi_{k}\|^{2}\big]=\sum_{j\neq k}\frac{\mathbb{E}|\Xi_{jk}|^{2}}{(\lambda_{k}-\lambda_{j})^{2}}.
\end{equation}
Esta equação conecta a variância nos incrementos dos autovetores de
$K$com a variância na projeção de $d\eta$, mostrando que as propriedades
do bulk não são indepentes, mas emergem da propagação das flutuações
de Itô do operador estrutural $K$ amplificadas pela quase degenerescência
espectral. Quando $|\lambda_{k}-\lambda_{j}|$ é pequeno, o denominador
amplifica o ruído, fenômeno intimamente relacionado à repulsão espectral
descrita na teoria de matrizes aleatórias \cite{mehta2004} Subespaços
onde os autovalores são quase degenerados tornam-se altamente sensíveis
às flutuações. Esse regime define o \emph{bulk espectral}. No subespaço
quase degenerado, a mistura estocástica torna-se dominante, os autovetores
passam a sofrer uma rotação browniana efetiva, tornando-se estatisticamente
isotrópicos \cite{bourgade2013}. Consequentemente, as projeções dos
retornos nesse subespaço tornam-se gaussianas por efeito central:

\begin{equation}
a_{k}=\langle r,\phi_{k}\rangle\quad\Longrightarrow\quad a_{k}\sim\mathcal{N}(0,\sigma_{b}^{2}).
\end{equation}

Sob as hipóteses de mistura estocástica forte no subespaço degenerado,
consistentes com o princípio de universalidade espectral \cite{tao2012},
ausência de estrutura privilegiada residual, e limite de grande dimensão,
a densidade espectral do bulk converge para o formato de Marchenko--Pastur
\cite{marchenko1967,baik2016}:

\begin{equation}
\rho_{MP}(\lambda)=\frac{1}{2\pi q\sigma_{b}^{2}\lambda}\sqrt{(\lambda_{+}-\lambda)(\lambda-\lambda_{-})},
\end{equation}

com $\lambda_{\pm}=\sigma_{b}^{2}(1\pm\sqrt{q})^{2}.$

\section{Processos de Itô em Espaço de Matrizes e o Teorema de Flutuação--Dissipação
na Métrica de Frobenius}

\subsection{Dinâmica Estocástica Matricial}

Considere um processo estocástico matricial 
\[
K(t)\in\mathbb{R}^{n\times n},
\]
cuja dinâmica é descrita por uma equação diferencial estocástica de
Itô da forma 
\begin{equation}
dK=A(K)\,dt+B(K)\,dW,\label{eq:Ito_matrix}
\end{equation}
onde $A(K)$ é o termo determinístico (drift), $B(K)$ é o operador
de acoplamento do ruído, $dW$ é um processo de Wiener matricial com
\[
\mathbb{E}[dW]=0,\qquad\mathbb{E}[dW\,dW^{\top}]=Q\,dt,
\]

onde $Q$ é uma matriz simétrica positiva definida. O espaço $\mathbb{R}^{n\times n}$
é um espaço vetorial real de dimensão $n^{2}$. Seu produto interno
natural é o produto de Frobenius, dado por 
\begin{equation}
A:B\;=\;\mathrm{Tr}(A^{\top}B)\;=\;\sum_{i,j}A_{ij}B_{ij}.\label{eq:frobenius}
\end{equation}
A norma induzida é a norma de Frobenius 
\[
\|A\|_{F}^{2}=A:A.
\]
Este produto interno é o análogo matricial do produto escalar Euclidiano.
Definimos o funcional escalar 
\begin{equation}
E(K)=K:K=\mathrm{Tr}(K^{\top}K),\label{eq:energy}
\end{equation}
que mede o equivalente a \emph{energia quadrática} total do sistema
no espaço das matrizes. Aplicando a regra de Itô ao funcional $E(K)$
usando o produto interno de Frobenius, obtemos

\[
dE=d(K:K),
\]
onde

\[
d(K:K)=(dK):K+K:(dK)+(dK):(dK).
\]
Como o produto é simétrico:

\[
(dK):K=K:(dK),
\]
obtemos:

\begin{equation}
dE=2K:dK+dK:dK,\label{eq:Ito_energy}
\end{equation}
que é a forma matricial da regra de Itô para a energia quadrática.
Substituindo a dinâmica \eqref{eq:Ito_matrix} em \eqref{eq:Ito_energy},
obtemos

\[
dE=2K:(A\,dt+B\,dW)+(A\,dt+B\,dW):(A\,dt+B\,dW).
\]
Pelas regras de Itô, eliminamos os termos com

\[
dt^{2}=0,\quad dt\,dW=0,\quad dW\,dt=0,
\]
ou seja, 

\[
dE=2K:A\,dt+2K:B\,dW+(B\,dW):(B\,dW).
\]
Tomando média $\mathbb{E}[]$ de $dE$, observamos que

\[
\mathbb{E}[K:B\,dW]=0,
\]
logo

\begin{equation}
\frac{d}{dt}\mathbb{E}[E]=2\mathbb{E}[K:A]+\mathbb{E}[BQB^{\top}:I].\label{eq:energy_balance}
\end{equation}
Que é o balanço de enerrgia para este processo. O termo $2\mathbb{E}[K:A]$
dissipa energia, enquanto $\mathbb{E}[BQB^{\top}:I]$ injeta energia
no sistema. No regime estacionário:

\[
\frac{d}{dt}\mathbb{E}[E]=0,
\]
logo

\begin{equation}
2\mathbb{E}[K:A]=-\mathbb{E}[BQB^{\top}:I].\label{eq:FDT_scalar}
\end{equation}
que é a forma escalar do teorema de flutuação--dissipação (TFD) na
métrica de Frobenius. Em sistemas reais pode ocorrer uma quase estacionariedade,

\[
\langle dE\rangle_{W}\approx0
\]
apenas em janelas temporais finitas $W$. Neste caso temos equilíbrio
local, mas possível deriva lenta global. A razão

\[
R=-\frac{\langle2K:dK\rangle_{W}}{\langle dK:dK\rangle_{W}}
\]
quantifica o desvio do equilíbrio.

\subsection{Distribuição dos Tempos Entre Mudanças de Regime}

Definimos uma mudança de regime espectral como o instante em que um
autovalor estrutural $\lambda_{k}(\tau)$ cruza o limite superior
do bulk de Marchenko--Pastur $\lambda_{+}$, isto é, $\lambda_{k}(\tau)=\lambda_{+}(\tau)$.
Buscamos aqui determinar a forma assintótica da distribuição de tempos
$g(\tau^{*})$, e sua conexão com a dinâmica de $K$, onde $\tau^{*}$é
o intervalo de tempo entre mudanças de regime. Considere a dinâmica
matricial geral 
\[
dK=A(K,\tau)\,d\tau+B(K,\tau)\,dW,
\]
onde $A$ e $B$ são operadores suficientemente regulares e $dW$
é um processo de Wiener matricial com momentos de segunda ordem finitos.
Aplicando cálculo de Itô para autovalores de operadores auto-adjuntos
\cite{dyson1962,anderson2010}, obtemos 
\begin{equation}
d\lambda_{k}=\langle\phi_{k},dK\phi_{k}\rangle+\sum_{j\neq k}\frac{|\langle\phi_{j},dK\phi_{k}\rangle|^{2}}{\lambda_{k}-\lambda_{j}}.\label{eq:lambda_ito}
\end{equation}
Definimos a variável 
\[
X(\tau)=\lambda_{k}(\tau)-\lambda_{+}(\tau),
\]
que é a distância entre um autovalor e o limiar do bulk. Com essa
definição, uma mudança de regime ocorre quando $X(\tau)=0$. Próximo
à borda do bulk, os autovalores satisfazem propriedades universais
da teoria de matrizes aleatórias \cite{mehta2004,tao2012}. Em particular,
o espaçamento espectral próximo à borda escala como $n^{-2/3}$. Já
a repulsão espectral implica que termos do tipo $(\lambda_{k}-\lambda_{j})^{-1}$
permanecem dominantes. Da Eq.~(\ref{eq:lambda_ito}), o termo singular
\[
\sum_{j\neq k}\frac{|\langle\phi_{j},dK\phi_{k}\rangle|^{2}}{\lambda_{k}-\lambda_{j}}
\]
gera amplificação estocástica quando $\lambda_{k}\to\lambda_{+}$.
No entanto, independentemente da forma específica de $A$ e $B$,
a dinâmica efetiva de $X(\tau)$ pode ser escrita localmente como
um processo de difusão unidimensional: 
\begin{equation}
dX=\mu(X,\tau)\,d\tau+\sigma(X,\tau)\,dW_{\tau},\label{eq:general_diffusion}
\end{equation}
onde $\mu(X,\tau)$ é um drift regular em $X$ e $\sigma(X,\tau)$
é contínua e estritamente positiva. A regularidade deste processo
segue da suavidade de $A$ e dos momentos finitos do ruído matricial.
Próximo ao cruzamento, temos $X\to0$, e pela continuidade de $\mu$,
existe expansão local tal que
\[
\mu(X,\tau)=a(\tau)X+o(X).
\]
Logo, no limite $X\to0$, o drift é de ordem linear enquanto a difusão
permanece finita $\sigma(0,\tau)>0$. Esse regime corresponde a um
ponto quase crítico, caracterizado por drift fraco, difusão dominante
e ausência de escala característica. Se $\tau^{*}$ é o tempo de primeira
passagem de $X(\tau)$ até zero, partindo de $X(0)>0$, para processos
de difusão unidimensionais da forma (\ref{eq:general_diffusion}),
resultados clássicos de teoria de processos estocásticos \cite{redner2001,majumdar1999}
mostram que, quando o drift é regular e a difusão é dominante próximo
ao ponto crítico, a distribuição assintótica do tempo de primeira
passagem satisfaz 
\begin{equation}
g(\tau^{*})\sim\tau^{*-2}.\label{eq:powerlaw}
\end{equation}
Esse resultado é universal e independe da forma detalhada de $\mu$
e $\sigma$, desde que $\mu(X)$ seja $O(X)$ quando $X\to0$, $\sigma(0)\neq0$,
e o processo $dX$ seja Markoviano. A degenerescência espectral $\lambda_{k}\to\lambda_{+}$
atua como ponto crítico dinâmico, quando o gap espectral 
\[
\Delta(\tau)=\min_{i\neq j}|\lambda_{i}-\lambda_{j}|
\]
se aproxima de zero, ocorre amplificação estocástica via termos $(\lambda_{k}-\lambda_{j})^{-1}$,
redução efetiva do drift restaurador e critical slowing down. Consequentemente,
as mudanças de regime não possuem escala temporal característica,
e a distribuição dos intervalos entre eventos obedece à lei de potência
(\ref{eq:powerlaw}). Esta lei assintótica emerge como consequência
universal da dinâmica estocástica dos autovalores próxima à borda
de Marchenko--Pastur, independentemente de hipóteses específicas
sobre a forma do drift $A(K)$ ou da estrutura do ruído $B(K)$. Mudanças
de regime espectrais são, portanto, fenômenos críticos associados
à degenerescência do operador de covariância temporal.

\section{Conexão com o modelo de Black and Schoules}

Nesta seção conectamos, de forma rigorosa, a decomposição espectral
proposta para os incrementos logarítmicos com o modelo de Black--Scholes.
O ponto central é que, no limite contínuo, a volatilidade que entra
na equação de precificação não é um parâmetro livre, mas uma quantidade
espectral determinada pelos autovalores do bulk de Marchenko--Pastur
do operador de covariância dos incrementos.

\subsection{O objeto espectral no limite contínuo}

No caso discreto, trabalhamos com

\[
P(t;\Delta t)=\ln\left(\frac{S(t)}{S(t-\Delta t)}\right),
\]
e a matriz de covariância construída é

\[
C_{tt'}=\langle P(t)P(t')\rangle.
\]
Quando $\Delta t\to0$,

\[
P(t;\Delta t)=\int_{t-\Delta t}^{t}d(\ln S)\longrightarrow d(\ln S_{t}).
\]
Logo, o operador contínuo correspondente é

\[
C(t,t')=\left\langle d(\ln S_{t})d(\ln S_{t'})\right\rangle ,
\]
isto é, a covariância dos incrementos do log-preço. Nesse limite,
os autovetores discretos $\phi_{k}(t)$ passam a representar padrões
de incrementos, que denotaremos por $d\phi_{k}(t)$. A expansão espectral
contínua assume, então, a forma

\begin{equation}
d(\ln S)=\sum_{k\le m}a_{k}d\phi_{k}(t)+\sigma dW_{t},\label{eq:expansao_incrementos}
\end{equation}
onde os modos $k\le m$ correspondem aos autovalores fora do bulk
MP, o termo $dW_{t}$ representa a contribuição combinada dos modos
do bulk.

\subsection{Identificação de $\mu$ e $\sigma$}

Tomando a esperança em (\ref{eq:expansao_incrementos}),

\[
\mathbb{E}[d(\ln S)]=\sum_{k\le m}a_{k}\mathbb{E}[d\phi_{k}].
\]
Como $\phi_{k}$ são processos suaves no tempo físico,

\[
\mathbb{E}[d\phi_{k}]=\frac{d\phi_{k}}{dt}dt.
\]
 Logo,

\begin{equation}
\mu=\sum_{k\le m}a_{k}\frac{d\phi_{k}}{dt}.\label{eq:mu}
\end{equation}
Por outro lado, a variância do termo ruidoso é determinada pelos autovalores
do bulk \cite{BunBouchaudPotters2017}

\begin{equation}
\sigma^{2}=\sum_{k>m}\lambda_{k}.\label{eq:sigma}
\end{equation}
Esta volatilidade pode ser interpretada como resultado da contribuição
de múltiplos modos rápidos não resolvidos\cite{Fouque2000}. Substituindo
em $d(\ln S)=\frac{dS}{S}$, obtemos a equação diferencial efetiva

\[
\frac{dS}{S}=\mu dt+\sigma dW_{t},
\]
que é exatamente a dinâmica assumida no modelo de Black--Scholes.

\subsection{Dependência espectral da volatilidade e das sensibilidades}

A equação de precificação de Black--Scholes,

\[
\frac{\partial C}{\partial t}+\frac{1}{2}\sigma^{2}S^{2}\frac{\partial^{2}C}{\partial S^{2}}+rS\frac{\partial C}{\partial S}-rC=0,
\]
depende da volatilidade apenas através de$\sigma^{2}$. Pela relação
(\ref{eq:sigma}), essa quantidade é uma soma espectral. Assim, qualquer
sensibilidade da opção à volatilidade é, na verdade, uma sensibilidade
aos autovalores do bulk MP.

A sensibilidade Vega clássica é

\[
\text{Vega}=\frac{\partial C}{\partial\sigma}.
\]
Aplicando a regra da cadeia,

\[
\frac{\partial C}{\partial\lambda_{k}}=\frac{\partial C}{\partial\sigma}\frac{\partial\sigma}{\partial\lambda_{k}}.
\]
 Como

\[
\sigma=\sqrt{\sum_{j>m}\lambda_{j}},
\]
temos

\[
\frac{\partial\sigma}{\partial\lambda_{k}}=\frac{1}{2\sigma}.
\]
Logo,

\[
\frac{\partial C}{\partial\lambda_{k}}=\frac{1}{2\sigma}\text{Vega}.
\]
A Vega pode, portanto, ser interpretada como a soma das sensibilidades
da opção a cada autovalor do bulk. Os autovalores $\lambda_{k}=\lambda_{k}(\tau)$
variam no tempo lento (regime de mercado). Assim, a sensibilidade
$\Theta$ pode ser escrita como

\[
\Theta=\frac{\partial C}{\partial t}+\frac{\partial C}{\partial\sigma}\frac{\partial\sigma}{\partial\tau},
\]
e como

\[
\frac{\partial\sigma}{\partial\tau}=\frac{1}{2\sigma}\sum_{k>m}\frac{\partial\lambda_{k}}{\partial\tau},
\]
 a $\Theta$ mede a sensibilidade da opção à evolução do espectro
MP. As outras sensibilidades do modelo de Black and Schoules, $\Gamma$,
$\Delta$ e $\rho$ possuem dependência indireta com $\sigma$, e
portanto com os autovalores do bulk. 

O drift $\mu$ depende apenas dos modos fora do bulk, pela equação
(\ref{eq:mu}). Qualquer variação nesses modos altera o valor da opção
através da dependência de $C$ em $\mu$. Surge, assim, uma sensibilidade
adicional:

\[
\frac{\partial C}{\partial\lambda_{k}},\quad k\le m,
\]

associada à rotação das primitivas estruturais. 

\section{Conexões com a Teoria Moderna de Portfólio: correlação entre primitivas}

A Teoria Moderna de Portfólio utiliza como objeto central a matriz
de covariância entre ativos,

\[
\Sigma_{ij}=\langle r_{i}(t),r_{j}(t)\rangle,
\]
onde $r_{i}(t)$ e $r_{j}(t)$ são os retornos de dois ativos distintos
ao longo de uma janela temporal. Essa matriz é interpretada como contendo
a estrutura de correlação relevante do mercado, e sua decomposição
espectral fornece os fatores principais de risco utilizados na construção
de portfólios eficientes. No entanto, à luz da decomposição espectral
temporal introduzida neste trabalho, essa covariância total pode ser
analisada de forma mais detalhada.

Para cada ativo $i$, dentro de uma janela, escrevemos

\[
r_{i}(t)=\sum_{k=1}^{m_{i}}a_{ik}\phi_{ik}(t)+\eta_{i}(t),
\]
onde ${\phi_{ik}(t)}$ são as primitivas temporais informativas, isto
é, as autofunções associadas a autovalores fora do bulk de Marchenko--Pastur,
$a_{ik}$ são os coeficientes dessa decomposição, $\eta_{i}(t)$ é
o resíduo puramente ruidoso, cuja estatística é descrita pela teoria
de matrizes aleatórias e satisfaz, aproximadamente, 
\[
\langle\eta_{i}(t),\eta_{i}(t')\rangle\approx\sigma_{i}^{2}\delta(t-t').
\]

Substituindo essa decomposição na definição da covariância de Markowitz,
obtemos

\[
\Sigma_{ij}=\Big\langle\sum_{k}a_{ik}\phi_{ik}(t)+\eta_{i}(t),\sum_{\ell}a_{j\ell}\phi_{j\ell}(t)+\eta_{j}(t)\Big\rangle.
\]

Expandindo os termos,

\[
\Sigma_{ij}=\sum_{k,\ell}a_{ik}a_{j\ell}\langle\phi_{ik},\phi_{j\ell}\rangle+\langle\eta_{i},\eta_{j}\rangle.
\]
 

A matriz de covariância observada é, portanto, a soma de duas contribuições
de naturezas distintas: a correlação estrutural, proveniente do compartilhamento
de primitivas temporais entre os ativos, e a correlação espúria, proveniente
da interação entre os resíduos ruidosos previstos pela teoria de Marchenko--Pastur.
O segundo termo não contém informação estrutural sobre a dinâmica
do mercado. Ele é consequência puramente estatística da estimação
finita da covariância e do ruído branco temporal que compõe os resíduos
$\eta_{i}(t)$. No entanto, esse termo está presente na matriz $\Sigma$
utilizada na prática na teoria de portfólio.

A parte estrutural da correlação pode ser isolada como

\[
\Sigma_{ij}^{\mathrm{estr}}=\sum_{k,\ell}a_{ik}a_{j\ell}\langle\phi_{ik},\phi_{j\ell}\rangle.
\]

Essa expressão revela que dois ativos são correlacionados não porque
seus retornos totais são semelhantes, mas porque compartilham primitivas
temporais semelhantes. Em outras palavras, a correlação entre ativos
é mediada pelos padrões temporais que ambos exibem, como:
\begin{itemize}
\item modos de tendência, 
\item modos de reversão à média,
\item modos associados a clusters de volatilidade, 
\item entre outros padrões estruturais capturados pelas autofunções fora
do bulk.
\end{itemize}
Sob essa perspectiva, a covariância total utilizada por Markowitz
mistura estrutura com ruído estatístico. Isso fornece uma explicação
conceitual para um fato empírico bem conhecido: matrizes de covariância
estimadas a partir de retornos são notoriamente instáveis fora da
amostra \cite{LedoitWolf2004,LedoitWolf2012}, pois parte significativa
de sua variabilidade provém do termo residual. A matriz de covariância
estrutural “limpa” adequada para a teoria de portfólio é, portanto,

\[
\widetilde{\Sigma}_{ij}=\Sigma_{ij}^{\mathrm{estr}},
\]
isto é, a covariância reconstruída apenas a partir das primitivas
informativas de cada ativo.

Uma consequência importante dessa decomposição é que o espaço relevante
para diversificação não é propriamente o espaço dos ativos, mas o
espaço das primitivas temporais compartilhadas por eles. A construção
de portfólios eficientes passa, assim, a depender não apenas da magnitude
das correlações observadas, mas da identificação dos padrões temporais
estruturais que os ativos têm em comum.

Essa reformulação fornece uma interpretação espectral mais refinada
da teoria moderna de portfólio, na qual a correlação entre ativos
é entendida como correlação entre suas primitivas temporais, enquanto
as contribuições provenientes do ruído estatístico são explicitamente
removidas da análise.

\section{Conexão com Modelos de Fatores Utilizados na Indústria}

Grandes investidores institucionais e fundos quantitativos constroem
carteiras com base em fatores empíricos como \textit{momentum}, \textit{value},
\textit{low-volatility}, \textit{quality}, entre outros \cite{Bouchaud2003,FamaFrench1993,FamaFrench2015,Carhart1997,AsnessMoskowitzPedersen2013,FrazziniPedersen2014}.
Muitos fatores são definidos como funções do histórico de preços e
retornos dos ativos, e são amplamente utilizados por gestores como
AQR, BlackRock, Two Sigma e Bridgewater.

Esses fatores não são definidos a partir de uma teoria dinâmica do
preço, mas como observáveis estatísticos do comportamento passado
das séries temporais. À luz da decomposição espectral temporal introduzida
neste trabalho, mostraremos que esses fatores podem ser expressos
explicitamente em termos dos autovalores $\lambda_{k}$, autovetores
temporais $\phi_{k}(t)$ e dos coeficientes $a_{ik}$ que conectam
cada ativo às primitivas temporais do regime.

\subsection{Momentum}

Definido como

\begin{equation}
\text{Mom}=\int_{t-T}^{t}P(s)\,ds,
\end{equation}

temos

\begin{equation}
\text{Mom}=\sum_{k}a_{k}M_{k},\quad M_{k}=\int\phi_{k}(s)\,ds.
\end{equation}

Momentum é, portanto, uma projeção linear dos coeficientes espectrais.
Entretanto, essa soma inclui os modos do bulk, que não trazem informação
real. Assim, definimos o momentum efetivo como sendo a soma somente
sobre os modos fora do bulk
\[
\text{Mom}^{eff}=\sum_{k=1}^{m}a_{k}M_{k}
\]


\subsection{Volatilidade e Low-Vol}

A variância temporal é

\begin{equation}
\sigma^{2}=\int P^{2}(t)\,dt\approx\sigma_{\text{modos principais}}^{2}+\sigma_{\text{bulk}}^{2}=\sum_{k=1}^{m}a_{k}^{2}+\sum_{k>m}\lambda_{k},
\end{equation}
assim, um ativo de baixa volatilidade é um ativo que minimiza 
\[
\sigma_{\text{modos principais}}^{2}=\sum_{k=1}^{m}a_{k}^{2},
\]
pois os modos do bulk não possuem informação real.

\section{Explicação Espectral da Fenomenologia Observada no Mercado}

Na introdução deste trabalho foram descritas diversas características
empíricas recorrentes observadas em séries temporais de preços financeiros.
Mostramos agora que essas características emergem naturalmente da
estrutura espectral do operador de covariância temporal $K$.

\subsection{Persistência de padrões dentro de regimes}

Dentro de um regime, o espectro informativo de $K$ é estável e as
primitivas temporais $\phi_{k}(t)$ permanecem aproximadamente invariantes.
Como os retornos são combinações lineares dessas primitivas,

\begin{equation}
P_{i}(t)=\sum_{k}a_{ik}\,\phi_{k}(t),
\end{equation}

os padrões observados nos preços persistem enquanto o espectro se
mantém estável.

\subsection{Mudanças abruptas de comportamento}

Transições de regime correspondem a rotações espectrais significativas
das primitivas $\phi_{k}(t)$ e a alterações no conjunto de autovalores
fora do bulk de Marchenko--Pastur. Isso implica que a base temporal
que descreve os retornos muda subitamente, produzindo a sensação empírica
de quebra estrutural no mercado.

\subsection{Predominância de ruído}

A maior parte do espectro de $K$ pertence ao bulk de Marchenko--Pastur,
o que implica que

\[
P_{i}(t)=\sum_{k\le m}a_{ik}\phi_{k}(t)+\sum_{k>m}a_{ik}\phi_{k}(t),
\]
onde o segundo termo é puramente ruidoso. Isso explica por que grande
parte da variabilidade observada nos preços não contém informação
estrutural.

\subsection{Traders técnicos e padrões gráficos}

Traders técnicos baseiam suas decisões na identificação de padrões
recorrentes nos gráficos de preços, como tendências, canais, suportes,
resistências e formações geométricas. Dentro do modelo espectral,
tais padrões não são coincidências visuais, mas consequências diretas
da persistência das primitivas temporais $\phi_{k}(t)$ dentro de
um regime. Como os retornos são combinações lineares dessas funções,

\begin{equation}
P_{i}(t)=\sum_{k}a_{ik}\,\phi_{k}(t),
\end{equation}
a repetição dos mesmos modos dominantes gera estruturas gráficas recorrentes
que podem ser exploradas empiricamente. A análise técnica pode, portanto,
ser interpretada como uma identificação heurística das primitivas
temporais do regime vigente.

\subsection{Algoritmos de alta frequência (HFT)}

Algoritmos HFT operam explorando microestruturas temporais de curtíssimo
prazo. No contexto espectral, essas microestruturas correspondem a
modos de alta frequência associados ao bulk de Marchenko--Pastur.

Como esses modos são numerosos e de natureza ruidosa, mas estatisticamente
estáveis dentro do regime, HFTs conseguem explorar pequenas assimetrias
estatísticas recorrentes nesses componentes espectrais de alta frequência,
sem depender dos modos informativos de baixa frequência.

\subsection{Clusters de volatilidade}

Um fenômeno amplamente documentado é a presença de clusters de volatilidade
\cite{Cont2001,Engle1982,Bollerslev1986}, onde períodos de alta variabilidade
são seguidos por outros períodos igualmente voláteis.

No modelo espectral, a volatilidade é dada por

\begin{equation}
\sigma_{i}^{2}=\sum_{k}a_{ik}^{2}.
\end{equation}

Como os coeficientes associados ao bulk são numerosos e dominantes,
pequenas rotações espectrais ou redistribuições de energia entre esses
modos produzem pequenas variações na soma total, gerando os clusters
de volatilidade observados empiricamente. Essa persistência não requer
memória explícita no processo de preços, mas surge naturalmente da
estrutura espectral do operador $K$. 

\subsection{Eficácia intermitente de fatores}

Fatores como momentum e baixa volatilidade funcionam enquanto as primitivas
informativas permanecem estáveis. Quando ocorre rotação espectral,
os funcionais que definiam esses fatores deixam de capturar as novas
combinações relevantes, explicando sua perda de eficácia.

\subsection{Volatilidade estrutural}

A decomposição

\[
\sigma_{i}^{2}=\sum_{k}a_{ik}^{2}
\]

mostra que a volatilidade é a soma das energias espectrais do ativo.
A parcela dominante proveniente do bulk explica a presença de uma
volatilidade estrutural aparentemente inevitável.

\bibliographystyle{plain}
\bibliography{referencias}

\end{document}
